\documentclass[12pt]{article}

% =========================
% Geometry & core packages
% =========================
\usepackage[a4paper,margin=0.8in]{geometry}
\usepackage{amsmath,amssymb,amsthm,mathtools}
\usepackage{bm}
\usepackage{array,booktabs}
\usepackage{enumitem}
\usepackage{microtype}
\usepackage{xcolor}
\usepackage{hyperref}
\usepackage{fancyhdr}
\usepackage{draftwatermark}
\usepackage{etoolbox}
\usepackage{titlesec}

% =========================
% Watermark
% =========================
\SetWatermarkText{Unofficial --- Refined Practice --- MH1101}
\SetWatermarkScale{1}
\SetWatermarkLightness{0.985}

% =========================
% Course metadata
% =========================
\newcommand{\CourseCode}{MH1101}
\newcommand{\CourseName}{Calculus II}
\newcommand{\DocType}{Refined Practice Test (With Solutions)}
\newcommand{\ExamMeta}{AY2025-26, Semester 2}
\newcommand{\ExamMetaShort}{Refined Practice}

% =========================
% Header/footer styles
% =========================
\fancypagestyle{main}{
\fancyhf{}
\fancyhead[L]{\CourseCode\ \CourseName\ --\ \ExamMetaShort}
\fancyhead[R]{Unofficial Practice}
\fancyfoot[C]{\thepage}
\renewcommand{\headrulewidth}{0.4pt}
\renewcommand{\footrulewidth}{0pt}
}

\fancypagestyle{plainfirst}{
\fancyhf{}
\renewcommand{\headrulewidth}{0pt}
\renewcommand{\footrulewidth}{0pt}
\fancyfoot[C]{\scriptsize\textcolor{gray}{\textit{For academic practice only. This is an original practice paper written in a similar format to the reference file, but with newly written, more difficult calculus questions and full solutions.}}}
}

% =========================
% Convenience macros
% =========================
\newcommand{\dd}{\,\mathrm{d}}
\newcommand{\ee}{\mathrm{e}}
\newcommand{\ii}{\mathrm{i}}
\newcommand{\R}{\mathbb{R}}
\newcommand{\N}{\mathbb{N}}
\DeclareMathOperator{\sech}{sech}

% =========================
% Overview block
% =========================
\newenvironment{overview}{
\par\bigskip
\begin{center}
\rule{0.92\linewidth}{0.6pt}\\[-0.2em]
\textbf{Overview}\\[-0.2em]
\rule{0.92\linewidth}{0.6pt}
\end{center}
\vspace{-0.3em}
\begin{itemize}[leftmargin=1.6em, itemsep=0.2em, topsep=0.2em]
}{
\end{itemize}
\par\bigskip
}

\setlist[enumerate]{itemsep=0.32em, topsep=0.32em}
\setlist[itemize]{itemsep=0.22em, topsep=0.22em}

\titleformat{\subsubsection}[block]
{\normalfont\large\bfseries}
{}
{0pt}
{}

\titleformat{\paragraph}[block]
{\normalfont\normalsize\bfseries}
{}
{0pt}
{}

\titlespacing*{\paragraph}
{0pt}
{1.2ex}
{0.6ex}

\setcounter{secnumdepth}{-1}
\setcounter{tocdepth}{4}

% =========================
% Title
% =========================
\title{
\Huge \CourseCode\ \CourseName \\[0.3em]
\Large \DocType \\[0.3em]
\ExamMeta
}
\date{\today}

% =========================
% Document begins
% =========================
\begin{document}

\maketitle
\thispagestyle{plainfirst}
\pagestyle{main}

\begin{overview}
\item \textbf{Time allowed:} 1 hour.
\item \textbf{Total marks:} 60.
\item \textbf{Difficulty:} Intentionally set above the level of the provided test papers and Tutorial 6.
\item \textbf{Coverage:} Riemann sums; FTC and Leibniz-rule differentiation; area and volume by both disk-washer and shell methods; improper integrals; nontrivial substitutions and integration by parts; numerical quadrature and error bounds; exactness of Simpson's Rule.
\end{overview}

\newpage
\tableofcontents
\newpage

% ============================================================
% Question 1
% ============================================================
\section{Question 1 \hfill (15 marks)}

\subsection{Problem}

\begin{enumerate}[label=(\alph*)]
\item \textbf{[5 marks]} Evaluate
\[
L=\lim_{n\to\infty}\sum_{i=1}^n \frac{2i}{n^2}\ln\!\left(1+\frac{i^2}{n^2}\right).
\]

\item \textbf{[5 marks]} Define
\[
F(x)=\cos x\int_{x^2}^{\,1+\sin x}(1+t^2)\dd t.
\]
Find \(F'(0)\).

\item \textbf{[5 marks]} Let \(b:[0,\infty)\to\R\) be differentiable on \((0,\infty)\), with \(b(x)\ge 0\) for all \(x\ge 0\) and \(b(0)=1\). Suppose that
\[
\frac{\dd}{\dd x}\int_0^{b(x)}\frac{2t}{1+t^2}\dd t=\frac{4x}{1+x^2}
\qquad (x>0).
\]
Find \(b(x)\).
\end{enumerate}

\subsection{Solution}

\subsubsection{(a) The limit}

\paragraph{Method 1: Direct Riemann-sum recognition}
Write
\[
\frac{2i}{n^2}\ln\!\left(1+\frac{i^2}{n^2}\right)
=
2\left(\frac{i}{n}\right)\ln\!\left(1+\left(\frac{i}{n}\right)^2\right)\frac{1}{n}.
\]
Hence, with
\[
f(x)=2x\ln(1+x^2),
\]
we obtain
\[
L=\int_0^1 2x\ln(1+x^2)\dd x.
\]
Now let \(u=1+x^2\). Then \(\dd u=2x\dd x\), so
\[
L=\int_1^2 \ln u \dd u
=
\left[u\ln u-u\right]_1^2
=
(2\ln2-2)-(-1).
\]
Therefore
\[
\boxed{L=2\ln2-1.}
\]

\paragraph{Method 2: Spot an antiderivative first}
Again,
\[
L=\int_0^1 2x\ln(1+x^2)\dd x.
\]
Observe that
\[
\frac{\dd}{\dd x}\Big((1+x^2)\ln(1+x^2)-x^2\Big)=2x\ln(1+x^2).
\]
Therefore
\[
L=\left[(1+x^2)\ln(1+x^2)-x^2\right]_0^1
=
(2\ln2-1)-0.
\]
Thus
\[
\boxed{L=2\ln2-1.}
\]

\bigskip

\subsubsection{(b) Differentiation under the integral sign}

\paragraph{Method 1: Product rule + Leibniz rule}
Let
\[
G(x)=\int_{x^2}^{\,1+\sin x}(1+t^2)\dd t.
\]
Then
\[
F(x)=\cos x\, G(x),
\]
so
\[
F'(x)=-\sin x\,G(x)+\cos x\,G'(x).
\]
By Leibniz's Rule,
\[
G'(x)=(1+(1+\sin x)^2)\cos x-(1+x^4)(2x).
\]
Hence
\[
F'(0)
=
-\sin 0\int_0^1(1+t^2)\dd t
+
\cos 0\left[(1+(1+\sin0)^2)\cos0-(1+0^4)(2\cdot0)\right].
\]
Since
\[
\int_0^1(1+t^2)\dd t=1+\frac13=\frac43,
\]
we get
\[
F'(0)=0+1\cdot(1+1)=2.
\]
Therefore
\[
\boxed{F'(0)=2.}
\]

\paragraph{Method 2: Write the integral explicitly first}
By the Fundamental Theorem of Calculus,
\[
\int (1+t^2)\dd t=t+\frac{t^3}{3}.
\]
Hence
\[
G(x)=\left[t+\frac{t^3}{3}\right]_{t=x^2}^{t=1+\sin x}
=
1+\sin x+\frac{(1+\sin x)^3}{3}-x^2-\frac{x^6}{3}.
\]
So
\[
F(x)=\cos x\left(1+\sin x+\frac{(1+\sin x)^3}{3}-x^2-\frac{x^6}{3}\right).
\]
Differentiate directly:
\[
F'(x)
=
-\sin x\left(1+\sin x+\frac{(1+\sin x)^3}{3}-x^2-\frac{x^6}{3}\right)
+\cos x\left(\cos x+(1+\sin x)^2\cos x-2x-2x^5\right).
\]
At \(x=0\),
\[
F'(0)=0+1\cdot(1+1)=2.
\]
Thus
\[
\boxed{F'(0)=2.}
\]

\bigskip

\subsubsection{(c) Solving the integral differential equation}

\paragraph{Method 1: FTC + chain rule}
By the Fundamental Theorem of Calculus and the chain rule,
\[
\frac{\dd}{\dd x}\int_0^{b(x)}\frac{2t}{1+t^2}\dd t
=
\frac{2b(x)b'(x)}{1+b(x)^2}.
\]
So the given equation becomes
\[
\frac{2b(x)b'(x)}{1+b(x)^2}=\frac{4x}{1+x^2}.
\]
But
\[
\frac{\dd}{\dd x}\ln(1+b(x)^2)=\frac{2b(x)b'(x)}{1+b(x)^2},
\qquad
\frac{\dd}{\dd x}\bigl(2\ln(1+x^2)\bigr)=\frac{4x}{1+x^2}.
\]
Hence
\[
\ln(1+b(x)^2)=2\ln(1+x^2)+C.
\]
Using \(b(0)=1\), we get
\[
\ln(1+1^2)=2\ln 1+C \quad\Rightarrow\quad \ln 2=C.
\]
Therefore
\[
\ln(1+b(x)^2)=2\ln(1+x^2)+\ln2
=
\ln\bigl(2(1+x^2)^2\bigr).
\]
Exponentiating,
\[
1+b(x)^2=2(1+x^2)^2.
\]
So
\[
b(x)^2=2(1+2x^2+x^4)-1=1+4x^2+2x^4.
\]
Since \(b(x)\ge 0\),
\[
\boxed{b(x)=\sqrt{1+4x^2+2x^4}.}
\]

\paragraph{Method 2: Identify the inside antiderivative first}
Since
\[
\int \frac{2t}{1+t^2}\dd t=\ln(1+t^2)+C,
\]
we have
\[
\int_0^{b(x)}\frac{2t}{1+t^2}\dd t
=
\ln(1+b(x)^2)-\ln 1
=
\ln(1+b(x)^2).
\]
Thus the hypothesis is simply
\[
\frac{\dd}{\dd x}\ln(1+b(x)^2)=\frac{4x}{1+x^2}.
\]
Integrating with respect to \(x\),
\[
\ln(1+b(x)^2)=2\ln(1+x^2)+C.
\]
Now use \(b(0)=1\) to get \(C=\ln 2\). Therefore
\[
1+b(x)^2=2(1+x^2)^2,
\]
and hence
\[
\boxed{b(x)=\sqrt{1+4x^2+2x^4}.}
\]

\newpage

% ============================================================
% Question 2
% ============================================================
\section{Question 2 \hfill (18 marks)}

\subsection{Problem}

Let \(R\) be the region enclosed by the curves
\[
y=x^{1/3},
\qquad
y=x^2,
\qquad
x\ge 0.
\]

\begin{enumerate}[label=(\alph*)]
\item \textbf{[4 marks]} Find the area of \(R\).

\item \textbf{[4 marks]} Use the disk-washer method to find the volume obtained when \(R\) is revolved about the \(x\)-axis.

\item \textbf{[4 marks]} Re-compute the volume in part (b) using cylindrical shells.

\item \textbf{[6 marks]} Let \(S\) be the region enclosed by the \(x\)-axis and the curve
\[
y=x^2-x^5
\qquad (0\le x\le 1).
\]
Find the volume generated when \(S\) is revolved about the \(y\)-axis, and explain clearly why the disk-washer method is not the efficient method here.
\end{enumerate}

\subsection{Solution}

First, solve the intersection points of \(x^{1/3}\) and \(x^2\):
\[
x^{1/3}=x^2
\quad\Longrightarrow\quad
x^{1/3}(1-x^{5/3})=0
\quad\Longrightarrow\quad
x=0 \text{ or } x=1.
\]
On \([0,1]\), we have \(x^{1/3}\ge x^2\).

\subsubsection{(a) Area of \(R\)}

\paragraph{Method 1: Integrate with respect to \(x\)}
\[
A=\int_0^1 \bigl(x^{1/3}-x^2\bigr)\dd x
=
\left[\frac{3}{4}x^{4/3}-\frac{x^3}{3}\right]_0^1
=
\frac34-\frac13
=
\frac{5}{12}.
\]
Thus
\[
\boxed{A=\frac{5}{12}.}
\]

\paragraph{Method 2: Integrate with respect to \(y\)}
From
\[
y=x^{1/3}\quad\Longrightarrow\quad x=y^3,
\qquad
y=x^2\quad\Longrightarrow\quad x=\sqrt{y},
\]
and \(0\le y\le 1\), the horizontal width is
\[
\sqrt y-y^3.
\]
So
\[
A=\int_0^1 (\sqrt y-y^3)\dd y
=
\left[\frac{2}{3}y^{3/2}-\frac{y^4}{4}\right]_0^1
=
\frac23-\frac14
=
\frac{5}{12}.
\]
Hence
\[
\boxed{A=\frac{5}{12}.}
\]

\bigskip

\subsubsection{(b) Volume about the \(x\)-axis by washers}

\paragraph{Method 1: Direct washer setup}
Outer radius:
\[
R(x)=x^{1/3}.
\]
Inner radius:
\[
r(x)=x^2.
\]
Therefore
\[
V
=
\pi\int_0^1\bigl(R(x)^2-r(x)^2\bigr)\dd x
=
\pi\int_0^1\bigl(x^{2/3}-x^4\bigr)\dd x.
\]
So
\[
V
=
\pi\left[\frac35 x^{5/3}-\frac{x^5}{5}\right]_0^1
=
\pi\left(\frac35-\frac15\right)
=
\frac{2\pi}{5}.
\]
Thus
\[
\boxed{V=\frac{2\pi}{5}.}
\]

\paragraph{Method 2: Reduce to a known shell answer}
Part (c) below computes the same solid by shells and gives
\[
V=2\pi\int_0^1 y(\sqrt y-y^3)\dd y
=
2\pi\left[\frac{2}{5}y^{5/2}-\frac{y^5}{5}\right]_0^1
=
2\pi\left(\frac25-\frac15\right)
=
\frac{2\pi}{5}.
\]
Hence the washer answer must be
\[
\boxed{V=\frac{2\pi}{5}.}
\]

\bigskip

\subsubsection{(c) Volume about the \(x\)-axis by shells}

\paragraph{Method 1: Horizontal shells}
Using \(y\)-integration, radius \(=y\), and shell height
\[
\sqrt y-y^3,
\]
we get
\[
V
=
2\pi\int_0^1 y(\sqrt y-y^3)\dd y
=
2\pi\int_0^1 \bigl(y^{3/2}-y^4\bigr)\dd y.
\]
Thus
\[
V
=
2\pi\left[\frac{2}{5}y^{5/2}-\frac{y^5}{5}\right]_0^1
=
2\pi\left(\frac25-\frac15\right)
=
\frac{2\pi}{5}.
\]
Therefore
\[
\boxed{V=\frac{2\pi}{5}.}
\]

\paragraph{Method 2: Compare with the washer integral}
From part (b),
\[
V=\pi\int_0^1 \bigl(x^{2/3}-x^4\bigr)\dd x=\frac{2\pi}{5}.
\]
Since both formulas compute the volume of the same solid, the shell method must also give
\[
\boxed{V=\frac{2\pi}{5}.}
\]

\bigskip

\subsubsection{(d) Volume about the \(y\)-axis for \(y=x^2-x^5\)}

\paragraph{Method 1: Cylindrical shells}
Since \(x^2-x^5=x^2(1-x^3)\ge 0\) on \([0,1]\), the shell radius is \(x\) and the shell height is
\[
x^2-x^5.
\]
Hence
\[
V
=
2\pi\int_0^1 x(x^2-x^5)\dd x
=
2\pi\int_0^1 (x^3-x^6)\dd x.
\]
Therefore
\[
V
=
2\pi\left[\frac{x^4}{4}-\frac{x^7}{7}\right]_0^1
=
2\pi\left(\frac14-\frac17\right)
=
2\pi\cdot\frac{3}{28}
=
\frac{3\pi}{14}.
\]
Thus
\[
\boxed{V=\frac{3\pi}{14}.}
\]

The shell method is efficient because it uses the graph directly as \(y=f(x)\).

\paragraph{Method 2: Why washers are not efficient}
If we try the disk-washer method about the \(y\)-axis, then we must write \(x\) as a function of \(y\) from
\[
y=x^2-x^5.
\]
Equivalently,
\[
x^5-x^2+y=0,
\]
which is a quintic equation in \(x\). In general, this does not lead to a useful elementary inverse formula for the inner and outer radii as functions of \(y\). Therefore the washer method is not the efficient choice here, whereas shells work immediately.

So the required volume is
\[
\boxed{V=\frac{3\pi}{14}.}
\]

\newpage

% ============================================================
% Question 3
% ============================================================
\section{Question 3 \hfill (15 marks)}

\subsection{Problem}

\begin{enumerate}[label=(\alph*)]
\item \textbf{[5 marks]} Evaluate the improper integral
\[
I=\int_0^\infty \frac{1}{e^x+e^{-x}+2}\dd x.
\]

\item \textbf{[4 marks]} Evaluate
\[
\int \cos(\ln x)\dd x
\qquad (x>0).
\]

\item \textbf{[6 marks]} Prove that for every \(a>0\),
\[
\int_0^a f''(x)(a-x)\dd x=f(a)-f(0)-af'(0),
\]
assuming \(f''\) is continuous on \([0,a]\). Then use the identity with \(f(x)=e^x\) to show that
\[
e^a\ge 1+a.
\]
\end{enumerate}

\subsection{Solution}

\subsubsection{(a) The improper integral}

\paragraph{Method 1: Exponential substitution}
Let
\[
u=e^x.
\]
Then \(\dd u=e^x\dd x=u\dd x\), so \(\dd x=\frac{1}{u}\dd u\). As \(x:0\to\infty\), we have \(u:1\to\infty\). Also
\[
e^x+e^{-x}+2=u+\frac{1}{u}+2=\frac{(u+1)^2}{u}.
\]
Hence
\[
I
=
\int_1^\infty \frac{1}{(u+1)^2/u}\cdot \frac{1}{u}\dd u
=
\int_1^\infty \frac{1}{(u+1)^2}\dd u.
\]
Thus
\[
I=\left[-\frac{1}{u+1}\right]_1^\infty
=0-\left(-\frac12\right)=\frac12.
\]
Therefore
\[
\boxed{I=\frac12.}
\]

\paragraph{Method 2: Hyperbolic simplification}
Observe that
\[
e^x+e^{-x}=2\cosh x.
\]
So
\[
e^x+e^{-x}+2=2(\cosh x+1)=4\cosh^2(x/2),
\]
using
\[
\cosh x+1=2\cosh^2(x/2).
\]
Therefore
\[
I=\int_0^\infty \frac{1}{4\cosh^2(x/2)}\dd x
=
\frac14\int_0^\infty \sech^2(x/2)\dd x.
\]
Now let \(u=x/2\). Then \(\dd x=2\dd u\), so
\[
I=\frac12\int_0^\infty \sech^2 u \dd u
=
\frac12\left[\tanh u\right]_0^\infty
=
\frac12(1-0)
=
\frac12.
\]
Hence
\[
\boxed{I=\frac12.}
\]

\bigskip

\subsubsection{(b) The integral \(\int \cos(\ln x)\dd x\)}

\paragraph{Method 1: Substitution \(u=\ln x\)}
Let
\[
u=\ln x.
\]
Then \(x=e^u\) and \(\dd x=e^u\dd u\). Hence
\[
\int \cos(\ln x)\dd x
=
\int e^u\cos u \dd u.
\]
Using the standard formula
\[
\int e^u\cos u\dd u=\frac{e^u}{2}(\sin u+\cos u)+C,
\]
we get
\[
\int \cos(\ln x)\dd x
=
\frac{x}{2}\bigl(\sin(\ln x)+\cos(\ln x)\bigr)+C.
\]
Therefore
\[
\boxed{\int \cos(\ln x)\dd x=\frac{x}{2}\bigl(\sin(\ln x)+\cos(\ln x)\bigr)+C.}
\]

\paragraph{Method 2: Integration by parts twice on the original variable}
Let
\[
I=\int \cos(\ln x)\dd x.
\]
Integrate by parts with
\[
u=\cos(\ln x),\qquad \dd v=\dd x.
\]
Then
\[
\dd u=-\frac{\sin(\ln x)}{x}\dd x,\qquad v=x.
\]
Hence
\[
I=x\cos(\ln x)+\int \sin(\ln x)\dd x.
\]
Let
\[
J=\int \sin(\ln x)\dd x.
\]
Again integrate by parts:
\[
u=\sin(\ln x),\qquad \dd v=\dd x.
\]
Then
\[
\dd u=\frac{\cos(\ln x)}{x}\dd x,\qquad v=x,
\]
so
\[
J=x\sin(\ln x)-\int \cos(\ln x)\dd x=x\sin(\ln x)-I.
\]
Substitute into the first equation:
\[
I=x\cos(\ln x)+x\sin(\ln x)-I.
\]
Thus
\[
2I=x\bigl(\sin(\ln x)+\cos(\ln x)\bigr),
\]
hence
\[
I=\frac{x}{2}\bigl(\sin(\ln x)+\cos(\ln x)\bigr)+C.
\]
Therefore
\[
\boxed{\int \cos(\ln x)\dd x=\frac{x}{2}\bigl(\sin(\ln x)+\cos(\ln x)\bigr)+C.}
\]

\bigskip

\subsubsection{(c) An identity and an application}

\paragraph{Method 1: Integration by parts}
Let
\[
J=\int_0^a f''(x)(a-x)\dd x.
\]
Integrate by parts with
\[
u=a-x,\qquad \dd v=f''(x)\dd x.
\]
Then
\[
\dd u=-\dd x,\qquad v=f'(x).
\]
So
\[
J=\left[f'(x)(a-x)\right]_0^a+\int_0^a f'(x)\dd x.
\]
Now
\[
\left[f'(x)(a-x)\right]_0^a=f'(a)\cdot 0-f'(0)\cdot a=-af'(0),
\]
and
\[
\int_0^a f'(x)\dd x=f(a)-f(0).
\]
Hence
\[
J=-af'(0)+f(a)-f(0).
\]
Therefore
\[
\boxed{\int_0^a f''(x)(a-x)\dd x=f(a)-f(0)-af'(0).}
\]

Now set \(f(x)=e^x\). Then \(f''(x)=e^x\), \(f(0)=1\), \(f'(0)=1\). So
\[
\int_0^a e^x(a-x)\dd x=e^a-1-a.
\]
For \(a>0\), we have \(e^x>0\) and \(a-x\ge 0\) on \([0,a]\), so the integral is nonnegative. Hence
\[
e^a-1-a\ge 0.
\]
Therefore
\[
\boxed{e^a\ge 1+a \qquad (a>0).}
\]
The same conclusion also holds for \(a=0\) by equality.

\paragraph{Method 2: Define an auxiliary function}
Define
\[
H(a)=f(a)-f(0)-af'(0)-\int_0^a f''(x)(a-x)\dd x.
\]
We show that \(H(a)\equiv 0\).

First,
\[
H(0)=f(0)-f(0)-0-\int_0^0 \cdots \dd x=0.
\]
Now differentiate with respect to \(a\). By Leibniz's Rule,
\[
\frac{\dd}{\dd a}\int_0^a f''(x)(a-x)\dd x
=
\int_0^a f''(x)\dd x,
\]
since the upper-endpoint term is \(f''(a)(a-a)=0\). Therefore
\[
H'(a)=f'(a)-f'(0)-\int_0^a f''(x)\dd x.
\]
By the Fundamental Theorem of Calculus,
\[
\int_0^a f''(x)\dd x=f'(a)-f'(0).
\]
Hence
\[
H'(a)=0.
\]
So \(H\) is constant, and since \(H(0)=0\), we get \(H(a)=0\) for all \(a\ge 0\). Thus
\[
\boxed{\int_0^a f''(x)(a-x)\dd x=f(a)-f(0)-af'(0).}
\]

Applying \(f(x)=e^x\) again yields
\[
e^a-1-a=\int_0^a e^x(a-x)\dd x\ge 0,
\]
since the integrand is nonnegative on \([0,a]\). Hence
\[
\boxed{e^a\ge 1+a.}
\]

\newpage

% ============================================================
% Question 4
% ============================================================
\section{Question 4 \hfill (12 marks)}

\subsection{Problem}

Let
\[
I=\int_0^1 e^x\dd x.
\]

\begin{enumerate}[label=(\alph*)]
\item \textbf{[4 marks]} Compute \(T_4\), \(M_4\), and \(S_4\), where \(T_n\), \(M_n\), and \(S_n\) are the composite Trapezoidal, Midpoint, and Simpson approximations, respectively.

\item \textbf{[4 marks]} Find values of \(n\) which guarantee that \(T_n\), \(M_n\), and \(S_n\) approximate \(I\) with error at most \(10^{-4}\).

\item \textbf{[4 marks]} Prove that Simpson's Rule with \(n=2\) is exact for every cubic polynomial
\[
p(x)=ax^3+bx^2+cx+d
\]
on every interval \([A,B]\).
\end{enumerate}

\subsection{Solution}

\subsubsection{(a) Computing \(T_4\), \(M_4\), and \(S_4\)}

Here \(h=\frac{1-0}{4}=\frac14\).

\paragraph{Method 1: Use the standard formulas directly}
For the Trapezoidal Rule,
\[
T_4=\frac{h}{2}\left[f(0)+2f\!\left(\frac14\right)+2f\!\left(\frac12\right)+2f\!\left(\frac34\right)+f(1)\right].
\]
Since \(f(x)=e^x\),
\[
T_4
=
\frac18\left(1+2e^{1/4}+2e^{1/2}+2e^{3/4}+e\right).
\]

For the Midpoint Rule, the midpoints are
\[
\frac18,\ \frac38,\ \frac58,\ \frac78,
\]
so
\[
M_4
=
\frac14\left(e^{1/8}+e^{3/8}+e^{5/8}+e^{7/8}\right).
\]

For Simpson's Rule,
\[
S_4=\frac{h}{3}\left[f(0)+4f\!\left(\frac14\right)+2f\!\left(\frac12\right)+4f\!\left(\frac34\right)+f(1)\right],
\]
hence
\[
S_4
=
\frac{1}{12}\left(1+4e^{1/4}+2e^{1/2}+4e^{3/4}+e\right).
\]

Therefore
\[
\boxed{
T_4=\frac18\left(1+2e^{1/4}+2e^{1/2}+2e^{3/4}+e\right),
}
\]
\[
\boxed{
M_4=\frac14\left(e^{1/8}+e^{3/8}+e^{5/8}+e^{7/8}\right),
}
\]
\[
\boxed{
S_4=\frac1{12}\left(1+4e^{1/4}+2e^{1/2}+4e^{3/4}+e\right).
}
\]

\paragraph{Method 2: Build each rule from subinterval averages}
For \(T_4\), average the endpoint values on each of the four subintervals and multiply by \(h=1/4\):
\[
T_4
=
\frac14\left[
\frac{1+e^{1/4}}{2}
+
\frac{e^{1/4}+e^{1/2}}{2}
+
\frac{e^{1/2}+e^{3/4}}{2}
+
\frac{e^{3/4}+e}{2}
\right],
\]
which simplifies to
\[
T_4=\frac18\left(1+2e^{1/4}+2e^{1/2}+2e^{3/4}+e\right).
\]

For \(M_4\), take the midpoint value on each subinterval:
\[
M_4
=
\frac14\left(e^{1/8}+e^{3/8}+e^{5/8}+e^{7/8}\right).
\]

For \(S_4\), apply Simpson on the two pairs of subintervals:
\[
S_4
=
\frac{1/2}{6}\left(1+4e^{1/4}+e^{1/2}\right)
+
\frac{1/2}{6}\left(e^{1/2}+4e^{3/4}+e\right),
\]
so
\[
S_4=\frac1{12}\left(1+4e^{1/4}+2e^{1/2}+4e^{3/4}+e\right).
\]
Hence the same formulas follow.

\bigskip

\subsubsection{(b) Error bounds}

\paragraph{Method 1: Apply the standard composite-rule bounds}
Since
\[
f(x)=e^x,
\qquad
f''(x)=e^x,
\qquad
f^{(4)}(x)=e^x,
\]
we may take
\[
K_2=\max_{0\le x\le 1}|f''(x)|=e,
\qquad
K_4=\max_{0\le x\le 1}|f^{(4)}(x)|=e.
\]

For the Trapezoidal Rule,
\[
|E_T|\le \frac{K_2(b-a)^3}{12n^2}
=
\frac{e}{12n^2}.
\]
We require
\[
\frac{e}{12n^2}\le 10^{-4},
\]
so
\[
n^2\ge \frac{e}{12\times 10^{-4}}\approx 2265.23.
\]
Hence it suffices to choose
\[
\boxed{n=48}
\]
(or any larger integer).

For the Midpoint Rule,
\[
|E_M|\le \frac{K_2(b-a)^3}{24n^2}
=
\frac{e}{24n^2}.
\]
We require
\[
\frac{e}{24n^2}\le 10^{-4},
\]
so
\[
n^2\ge \frac{e}{24\times 10^{-4}}\approx 1132.62.
\]
Hence it suffices to choose
\[
\boxed{n=34}.
\]

For Simpson's Rule,
\[
|E_S|\le \frac{K_4(b-a)^5}{180n^4}
=
\frac{e}{180n^4},
\]
with \(n\) even. We require
\[
\frac{e}{180n^4}\le 10^{-4},
\]
so
\[
n^4\ge \frac{e}{180\times 10^{-4}}\approx 151.02.
\]
Thus
\[
n\ge 3.50\ldots
\]
and the smallest even choice is
\[
\boxed{n=4}.
\]

\paragraph{Method 2: Compare orders before computing}
For \(e^x\) on \([0,1]\), all second and fourth derivatives are bounded by \(e\). Therefore:
\[
|E_T|=O(n^{-2}),\qquad |E_M|=O(n^{-2}),\qquad |E_S|=O(n^{-4}).
\]
So Simpson should require a much smaller \(n\) than the other two.

Using the exact constants:
\[
|E_T|\le \frac{e}{12n^2},\qquad
|E_M|\le \frac{e}{24n^2},\qquad
|E_S|\le \frac{e}{180n^4},
\]
we solve each inequality against \(10^{-4}\) exactly as above to obtain
\[
\boxed{n=48 \text{ for } T_n,\qquad n=34 \text{ for } M_n,\qquad n=4 \text{ for } S_n.}
\]

\bigskip

\subsubsection{(c) Simpson exactness for cubics}

Recall that Simpson's Rule on \([A,B]\) with \(n=2\) uses
\[
m=\frac{A+B}{2},
\qquad
S_2=\frac{B-A}{6}\bigl(p(A)+4p(m)+p(B)\bigr).
\]

\paragraph{Method 1: Basis-polynomial verification}
It suffices to prove exactness for the basis
\[
1,\ x,\ x^2,\ x^3,
\]
since Simpson's Rule and the definite integral are both linear in the function.

For \(p(x)=1\),
\[
\int_A^B 1\dd x=B-A,
\qquad
S_2=\frac{B-A}{6}(1+4+1)=B-A.
\]

For \(p(x)=x\),
\[
S_2=\frac{B-A}{6}\left(A+4\frac{A+B}{2}+B\right)
=
\frac{B-A}{6}(3A+3B)
=
\frac{B^2-A^2}{2}
=
\int_A^B x\dd x.
\]

For \(p(x)=x^2\),
\[
S_2
=
\frac{B-A}{6}\left(A^2+4\left(\frac{A+B}{2}\right)^2+B^2\right)
=
\frac{B-A}{6}\left(A^2+(A+B)^2+B^2\right).
\]
So
\[
S_2
=
\frac{B-A}{6}\bigl(2A^2+2AB+2B^2\bigr)
=
\frac{B-A}{3}(A^2+AB+B^2)
=
\frac{B^3-A^3}{3}
=
\int_A^B x^2\dd x.
\]

For \(p(x)=x^3\),
\[
S_2
=
\frac{B-A}{6}\left(A^3+4\left(\frac{A+B}{2}\right)^3+B^3\right)
=
\frac{B-A}{6}\left(A^3+\frac{(A+B)^3}{2}+B^3\right).
\]
Expand:
\[
A^3+\frac{(A+B)^3}{2}+B^3
=
A^3+\frac{A^3+3A^2B+3AB^2+B^3}{2}+B^3
=
\frac{3}{2}(A^3+A^2B+AB^2+B^3).
\]
Hence
\[
S_2
=
\frac{B-A}{4}(A^3+A^2B+AB^2+B^3).
\]
But
\[
B^4-A^4=(B-A)(B^3+B^2A+BA^2+A^3),
\]
so
\[
S_2=\frac{B^4-A^4}{4}=\int_A^B x^3\dd x.
\]

Therefore Simpson's Rule is exact for every cubic polynomial. Thus
\[
\boxed{S_2=\int_A^B p(x)\dd x \quad \text{for every cubic } p.}
\]

\paragraph{Method 2: Shift to the midpoint}
Let
\[
m=\frac{A+B}{2},
\qquad
h=\frac{B-A}{2},
\]
so that \(A=m-h\) and \(B=m+h\). Any cubic polynomial can be written in the shifted form
\[
p(x)=\alpha+\beta(x-m)+\gamma(x-m)^2+\delta(x-m)^3.
\]

Now Simpson's Rule on \([m-h,m+h]\) gives
\[
S_2=\frac{2h}{6}\bigl(p(m-h)+4p(m)+p(m+h)\bigr).
\]
Substitute the shifted form:

\[
p(m)=\alpha,
\]
\[
p(m-h)=\alpha-\beta h+\gamma h^2-\delta h^3,
\qquad
p(m+h)=\alpha+\beta h+\gamma h^2+\delta h^3.
\]
Hence
\[
p(m-h)+4p(m)+p(m+h)=6\alpha+2\gamma h^2.
\]
Therefore
\[
S_2=\frac{2h}{6}(6\alpha+2\gamma h^2)=2\alpha h+\frac{2}{3}\gamma h^3.
\]

Now compute the exact integral:
\[
\int_{m-h}^{m+h} p(x)\dd x
=
\int_{-h}^{h} (\alpha+\beta u+\gamma u^2+\delta u^3)\dd u
\qquad (u=x-m).
\]
The odd terms integrate to zero, so
\[
\int_{m-h}^{m+h} p(x)\dd x
=
\int_{-h}^{h}\alpha\dd u+\int_{-h}^{h}\gamma u^2\dd u
=
2\alpha h+\gamma\cdot \frac{2h^3}{3}.
\]
Hence
\[
\int_{m-h}^{m+h} p(x)\dd x=2\alpha h+\frac{2}{3}\gamma h^3=S_2.
\]
Therefore
\[
\boxed{S_2=\int_A^B p(x)\dd x \text{ for every cubic polynomial }p.}
\]

\newpage

% ============================================================
% Brief mark guidance
% ============================================================
\section{Brief Mark Guidance}

\subsection{Question 1 \hfill (15 marks)}
\begin{itemize}[leftmargin=1.6em]
\item \textbf{(a) (5 marks)} Correct Riemann-sum identification; correct integral; correct evaluation.
\item \textbf{(b) (5 marks)} Correct use of product rule and Leibniz rule or equivalent explicit antiderivative method; correct evaluation at \(x=0\).
\item \textbf{(c) (5 marks)} Correct FTC/chain-rule setup; correct integration and use of \(b(0)=1\); correct nonnegative branch.
\end{itemize}

\subsection{Question 2 \hfill (18 marks)}
\begin{itemize}[leftmargin=1.6em]
\item \textbf{(a) (4 marks)} Correct intersection points and area integral.
\item \textbf{(b) (4 marks)} Correct washer radii and final volume.
\item \textbf{(c) (4 marks)} Correct shell radius/height and final volume.
\item \textbf{(d) (6 marks)} Correct shell integral; correct final answer; clear explanation of why the washer method is not efficient.
\end{itemize}

\subsection{Question 3 \hfill (15 marks)}
\begin{itemize}[leftmargin=1.6em]
\item \textbf{(a) (5 marks)} Correct improper-integral setup and limit handling.
\item \textbf{(b) (4 marks)} Correct substitution or repeated integration by parts; correct antiderivative.
\item \textbf{(c) (6 marks)} Correct proof of the identity and correct application to \(f(x)=e^x\).
\end{itemize}

\subsection{Question 4 \hfill (12 marks)}
\begin{itemize}[leftmargin=1.6em]
\item \textbf{(a) (4 marks)} Correct formulas and correct substitution of nodes.
\item \textbf{(b) (4 marks)} Correct error bounds and smallest valid \(n\).
\item \textbf{(c) (4 marks)} Valid proof that Simpson's Rule with \(n=2\) is exact for every cubic.
\end{itemize}

\end{document}
